\section{Introduction}
\label{sec:introduction}

\vannak{} is an experimental operational knowledge plane for data pipelines and
metadata systems. It is not a generic log collector. Its purpose is to join
runtime process state, data provenance, metadata meaning, and cluster ownership
into a single low-latency operational view.

\vannak{} combines four established lines of work:

\begin{description}
\item[\durga{}] process monitoring: what pipeline process or activity is
  running, completed, failed, escalated, or cancelled.
\item[\ipto{}] metadata management: what datasets, schemas, attributes, lineage,
  classifications, contracts, and metadata objects mean.
\item[\sitas{}] shard-local runtime structure: hot in-memory indexes, explicit
  message passing, owned snapshots, and CPU/NUMA-aware placement.
\item[\raft{}] cluster control: agreed ownership, placement maps, checkpoints,
  and sealed-segment manifests (through the \graft{} implementation).
\end{description}

\subsection{The Operational Knowledge Plane}

An operational knowledge plane connects information that monitoring tools,
metadata catalogs, and workflow engines normally hold separately. A monitor can
tell you an activity failed. A catalog can tell you what a dataset means. A
workflow engine knows the process definition. None of them, on their own, can
answer a cross-cutting operational question quickly.

\vannak{} keeps recent operational state in memory, writes raw event history
and metadata outbox entries to append-only segments, writes durable metadata to
\ipto{}, and uses \raft{} for cluster coordination.

\subsection{The Two Event Planes}
\label{sec:intro-two-planes}

\vannak{} deliberately separates two related but distinct event streams.

\subsubsection{Process Events}

Process events are \durga{} monitor lifecycle events. They describe what the
pipeline process is doing. The Rust compatibility model mirrors \durga{}'s
canonical \texttt{org.gautelis.durga.ProcessEvent} shape and is represented by
\texttt{PipelineEvent} after ingestion. \vannak{} adds ingestion context such
as source identity, tenant, environment, metadata references, and causal parent.

\subsubsection{Data-Individual Metadata Events}

Data-individual metadata events describe what happened to one flowing data item.
They are provenance facts, not disposable monitoring telemetry. They are
represented by \texttt{DataIndividualMetadataEvent} and carry passive metadata
(received time, source, format, checksum) and active metadata (masking,
transformation, validation, enrichment, field changes).

\begin{vnote}
A \durga{} lifecycle event can say that an activity ran. A data-individual
metadata event says what happened to a specific data item during or around that
activity. The two planes complement each other; they are never merged.
\end{vnote}

\subsection{Operational Questions}

\vannak{} is designed to answer questions that need process state, provenance,
metadata meaning, and cluster ownership at the same time:

\begin{itemize}
\item Which pipelines are unhealthy right now?
\item Which datasets or metadata objects are affected by a failed activity?
\item What passive and active metadata has been captured for this specific data
  individual?
\item Which \ipto{} repository instance owns the durable metadata for this data
  item?
\item What changed before an incident started?
\item Which process instance last produced or modified a dataset?
\item Which cluster node owns the hot state, checkpoints, and sealed event
  segments for a pipeline or placement shard?
\end{itemize}

\subsection{The Thesis}

The first thesis is that process events become much more valuable when they are
joined with metadata meaning and kept in a low-latency shard-local index. The
second, equally important thesis is that metadata about individual flowing data
items must be captured durably as provenance facts, not treated as disposable
monitoring telemetry.

\subsection{Design Principles}

\vannak{} grows under a small set of principles that should be preserved unless
a task explicitly revises the architecture:

\begin{enumerate}
\item Keep high-volume event ingest out of the consensus hot path.
\item Use \raft{} for cluster control state, ownership, manifests, checkpoints,
  and metadata-version decisions.
\item Use \sitas{} for shard-local hot indexes, bounded ingest, and explicit
  cross-shard transfer.
\item Treat \durga{} process events as semantically meaningful state
  transitions, not unstructured logs.
\item Use \ipto{} metadata to connect runtime events to datasets, schemas,
  lineage, ownership, classifications, contracts, and versions.
\item Treat data-individual metadata capture as a durable provenance plane.
\item Route writes to \ipto{} by domain metadata placement, namely
  \texttt{DataIndividualShardId}; do not confuse it with \sitas{} executor
  shard ids.
\item Do not acknowledge durable metadata capture until the event is safely
  persisted to an outbox/segment or written idempotently to \ipto{}.
\item Keep ordinary application state shard-local.
\item Make cross-shard movement explicit through typed messages and handles.
\item Observability APIs return owned snapshots.
\end{enumerate}

\subsection{Crate Structure}

\vannak{} is a single Rust crate (\texttt{vannak}, edition 2024) with a
dependency-free domain core. Integrations attach at the module boundaries
rather than changing the core event model.

\begin{table}[htbp]
\centering
\begin{tabular}{ll}
\hline
Module & Responsibility \\
\hline
\texttt{ingest} & typed process-event envelope and validation \\
\texttt{process} & \durga{} process model and state reducer \\
\texttt{data} & data-individual identity, metadata events, provenance facts \\
\texttt{metadata} & typed metadata reference model \\
\texttt{durga} & \durga{} monitor compatibility types \\
\texttt{index} & dependency-free shard-local hot index \\
\texttt{query} & query request and result types \\
\texttt{ipto} & \ipto{} placement, mapping, payload, and outbox model \\
\texttt{ipto\_adapter} & \ipto{} writer backed by \texttt{ipto\_rust} \\
\texttt{storage} & append-only segment writer/reader and manifests \\
\texttt{cluster} & \raft{}-facing control state and placement \\
\texttt{raft\_sm} & \graft{} state machine adapter \\
\texttt{observability} & owned snapshot types \\
\texttt{runtime} & dependency-free shard-local runtime, bounded ingest, and
\sitas{} boundary \\
\texttt{sitas\_runtime} & \sitas{} executor/mailbox adapter
(\texttt{sitas-runtime}) \\
\hline
\end{tabular}
\caption{Crate module layout and responsibilities.}
\end{table}

\subsection{Current Status and Maturity}

This repository contains the dependency-free Rust core types and reducers plus
optional adapters behind feature flags. The core is implemented and tested.

The \texttt{vannak} binary (built with \texttt{node} + \texttt{daemon}) runs
the Raft control plane, the daemon service plane (health/query HTTP server,
background writer, segment discovery), and optionally the Kafka consumer in a
single supervised process. It replaces the earlier split between
\texttt{vannak-node} and \texttt{vannak-kafka-ingest} for deployments that
prefer co-located supervision. The standalone binaries remain available for
deployments that separate the control and data planes. \\

Implemented today: the \durga{}-compatible process-event mirror, the
process-state reducer, typed metadata references, the single-node hot index,
typed process, status, time-range, and impact queries, the dependency-free
shard-local runtime model (\texttt{ShardLocalRuntime}) with bounded ingest
queues (\texttt{BoundedIngestRuntime}), the feature-gated \sitas{}
executor/mailbox adapter (\texttt{SitasShardRuntime}), the Kafka-shaped
topic/partition/offset ingest boundary (\texttt{kafka\_ingest}), the
feature-gated \texttt{rdkafka} consumer runner (\texttt{vannak-kafka-ingest}),
the
data-individual provenance model and local provenance lookup index
(\texttt{DataProvenanceIndex}), consistent-hash ring placement with range
overrides, the \raft{}-replicable placement map with epoch-based query fallback,
field-to-attribute mapping, the idempotent write payload, the in-memory and
segment-backed metadata outbox, append-only segment storage, the typed
process-event journal (\texttt{ProcessEventJournal}), a single-node service
boundary (\texttt{VannakService}) that wires process ingest, durable process
event journaling, durable metadata capture, target-aware outbox draining,
writer-lease checks, checkpoint lookup, recovery replay, and owned snapshots, a
concrete \ipto{} writer
(\texttt{IptoRepoWriter}) with PROV-O relation-unit emission
(\texttt{wasGeneratedBy}, \texttt{used}, \texttt{wasDerivedFrom}), a
\graft{} state machine adapter (\texttt{ClusterStateMachine}), a daemon
runtime (\texttt{daemon.rs}) with segment discovery, background writer
scheduling, HTTP health/query endpoint, cluster query fanout, and automated
placement-change rebalancing, and the \texttt{vannak} supervisor binary that
co-locates Raft, Kafka, and the daemon service plane in one process.

\subsection{Feature Flags}

\vannak{} keeps optional dependencies behind Cargo feature flags so the
dependency-free core stays usable on its own.

\begin{table}[htbp]
\centering
\begin{tabular}{ll}
\hline
Feature & Effect \\
\hline
(none) & dependency-free domain core only \\
\texttt{ipto-writer} & enables \texttt{ipto\_adapter} backed by \texttt{ipto\_rust} \\
\texttt{raft} & enables \texttt{raft\_sm} with the \graft{} state machine \\
\texttt{node} & enables the \texttt{vannak-node} binary (implies \texttt{raft}) \\
\texttt{daemon} & daemon runtime: health/query HTTP, background writer, segment discovery \\
\texttt{sitas-runtime} & \texttt{SitasShardRuntime} adapter, \texttt{kafka\_ingest} boundary \\
\texttt{kafka-client} & \texttt{vannak-kafka-ingest} binary (implies \texttt{sitas-runtime}) \\
\hline
\end{tabular}
\caption{\vannak{} Cargo feature flags.}
\end{table}

\endinput
